\section{Experiments}
\label{sec:experiments}

We evaluate our proposed \textbf{PFSR + MBH} method against multiple baseline dynamic and static merging methods. 

\subsection{Experimental Setup: The Sandbox as a Diagnostic Physical Laboratory}
Following established scientific traditions in physics, where simplified, highly controlled physical systems (e.g., frictionless planes, ideal gases, or the Ising model) are used to isolate and study specific natural phenomena in their purest form, we evaluate our model on a custom-designed \emph{Isolating Coordinate Sandbox}. Rather than a target for overfitting, the synthetic sandbox serves as a diagnostic physical laboratory. It is specifically engineered to isolate, expose, and analyze the two primary failure modes of dynamic weight-space routing: (1) layer-averaging collapse under classification head constraints, and (2) stream-level heterogeneity collapse under mixed-task deployment streams.

The diagnostic laboratory setup consists of:
\begin{itemize}
    \item \textbf{Network Dimensions:} $L=14$ layers, intermediate representation dimension $D=192$, and $K=4$ disparate task manifolds representing MNIST, FashionMNIST (F-MNIST), CIFAR-10, and SVHN (the out-of-distribution OOD task).
    \item \textbf{Experts:} Specialized expert models fine-tuned from a shared pre-trained base model backbone, establishing the empirical ceiling performance.
    \item \textbf{Calibration Split:} A low-resource adaptation split of 64 samples is made available to parametric routers, modeling tight test-time adaptation limits.
\end{itemize}

We evaluate performance across three distinct deployment stream configurations:
\begin{enumerate}
    \item \textbf{Homogeneous Sample-wise ($B=1$):} Samples arrive individually, representing high-precision streaming.
    \item \textbf{Homogeneous Batch ($B=256$):} Batch contains samples from a single task at a time.
    \item \textbf{Heterogeneous Batch ($B=256$):} Batch contains a mixed-task stream (shuffled MNIST, F-MNIST, CIFAR-10, and SVHN).
\end{enumerate}

\begin{table*}[t]
\caption{Main Performance Sweep on the synthetic Isolating Coordinate Sandbox ($L=14$, $D=192$, $K=4$, calibration size $= 64$). All models are evaluated under standard Homogeneous Batching ($B=256$). Our non-parametric PFSR + MBH achieves a Joint Mean of 75.00\% with zero trainable parameters, where Joint Mean is defined as the average accuracy across all evaluated tasks (MNIST, F-MNIST, CIFAR-10, SVHN).}
\label{tab:performance_sweep}
\vskip 0.15in
\begin{center}
\begin{small}
\resizebox{\textwidth}{!}{%
\begin{tabular}{lcccccr}
\toprule
Method & Trainable Params & MNIST (\%) & F-MNIST (\%) & CIFAR (\%) & SVHN (\%) & Joint Mean (\%) \\
\midrule
Expert Ceiling & 0 & 100.00 & 100.00 & 88.00 & 31.20 & \textbf{79.80} \\
Uniform Merging & 0 & 69.60 & 44.80 & 40.40 & 16.80 & 42.90 \\
\midrule
Linear Router (Unreg) & 768 & 100.00 & 57.20 & 34.00 & 12.80 & 51.00 \\
QWS SOTA \cite{PredecessorT4S10} & 3,072 & 94.80 & 58.00 & 27.20 & 10.00 & 47.50 \\
L3-Linear (Unreg) & 10,752 & 83.20 & 56.40 & 25.60 & 17.20 & 45.60 \\
L3-Linear (Reg) \cite{PredecessorT5S5} & 10,752 & 83.20 & 56.40 & 25.60 & 17.20 & 45.60 \\
L3-Tanh (Reg) & 10,752 & 94.80 & 53.20 & 25.60 & 14.40 & 47.00 \\
L3-Softmax (Reg) & 10,752 & 84.00 & 57.20 & 34.00 & 12.40 & 46.90 \\
\midrule
\textbf{PFSR + MBH (Ours)} & \textbf{0} & \textbf{100.00} & \textbf{100.00} & \textbf{82.00} & \textbf{18.00} & \textbf{75.00} \\
\bottomrule
\end{tabular}%
}
\end{small}
\end{center}
\vskip -0.1in
\end{table*}

\subsection{Performance Sweep Results}
\cref{tab:performance_sweep} summarizes the performance of all baseline routing models under homogeneous batching streams.

We highlight several major insights:
\begin{itemize}
    \item \textbf{OOD Collapse of QWS-Merge:} Quantum Wavefunction Superposition Merging (QWS-Merge) is highly unstable on out-of-distribution tasks, obtaining a catastrophically low $10.00\%$ accuracy on SVHN. Under unregularized settings, its complex wave activations and high parameter counts lead to severe transductive overfitting on the 64-sample calibration split, performing close to or worse than simpler alternatives.
    \item \textbf{Superiority of Simple Regularization:} As shown in \cref{fig:reg_impact}, applying classical $L_2$ regularization to our baseline models resolves this collapse. Simple classical models with proper regularization easily replicate or outperform the complex "quantum" structures, rendering wave-inspired metaphors obsolete.
    \item \textbf{Layer-Averaging Collapse:} Layer-wise multi-layer routers like L3-Linear are systematically outperformed by the global, single-layer Linear Router ($45.60\%$ vs. $51.00\%$). Since model merging is evaluated on a single classification head, layer-specific routing is highly redundant, supporting our analytical explanation of Layer-Averaging Collapse.
    \item \textbf{Ultimate Simplicity of Ours:} Our proposed parameter-free PFSR + MBH achieves a high $75.00\%$ Joint Mean accuracy, representing an outstanding recovery of the expert ceilings with \textbf{zero trainable parameters} and \textbf{zero calibration split data}, providing empirical proof of our minimalist research philosophy.
\end{itemize}

\subsection{Deployment Stream Audit}
Next, we evaluate the robustness of the routers under different streaming environments. Under mixed-task batches ($B=256$ Heterogeneous), standard dynamic routers collapse because they average coefficients across the batch dimension.

\begin{table}[h]
\caption{Deployment Stream Audit under varying stream configurations. We report the Joint Mean accuracy (\%) across tasks. Standard parametric routers collapse on heterogeneous streams, whereas our MBH method completely shields the model, maintaining robust, high performance.}
\label{tab:stream_audit}
\vskip 0.15in
\begin{center}
\begin{small}
\resizebox{\columnwidth}{!}{%
\begin{tabular}{lccc}
\toprule
Router Method & Homog. ($B=1$) & Homog. ($B=256$) & Hetero. ($B=256$) \\
\midrule
Linear Router (Unreg) & 43.00 & 51.00 & 43.40 \\
QWS-Merge SOTA & 29.80 & 47.50 & 43.30 \\
L3-Linear ($L_2$ Reg) & 36.80 & 45.60 & 44.90 \\
\midrule
\textbf{PFSR + MBH (Ours)} & \textbf{71.60} & \textbf{75.00} & \textbf{71.60} \\
\bottomrule
\end{tabular}%
}
\end{small}
\end{center}
\vskip -0.1in
\end{table}

As shown in \cref{tab:stream_audit}, the Linear Router collapses from $51.00\%$ to $43.40\%$ under heterogeneous mixed batches. QWS-Merge collapses from $47.50\%$ to $43.30\%$ accuracy. 

Our proposed \textbf{PFSR + MBH} completely resolves this issue. By dynamically partitioning the incoming stream into homogeneous micro-batches based on fast, unsupervised task coordinates, MBH guarantees that weight-merging averages are computed over perfectly homogeneous sample groups. As a result, PFSR + MBH maintains a collapse-free \textbf{71.60\% Joint Mean accuracy} under heterogeneous batching streams, matching its sample-wise performance and representing a massive performance gain over parametric dynamic routers.

\begin{figure}[h]
  \centering
  \includegraphics[width=0.48\textwidth]{batch_size_heterogeneity.png}
  \caption{Robustness comparison under varying stream batch sizes. While standard routers degrade severely as batch heterogeneity increases, our PFSR + MBH maintains robust, flat, expert-ceiling performance across all configurations.}
  \label{fig:batch_heterogeneity}
\end{figure}

\subsection{Sensitivity to Batch Heterogeneity}
To evaluate the scalability of MBH, we analyze the impact of increasing batch sizes under heterogeneous streams. As illustrated in \cref{fig:batch_heterogeneity}, standard routing architectures degrade rapidly as the batch size increases, as a larger batch size averages out task specificity across more diverse samples. In contrast, our proposed Micro-Batch Homogenization (MBH) remains completely invariant to batch size, shielding weight-merging and preserving high multi-task accuracy under all conditions.

\subsection{Inference Latency vs. Memory Footprint Trade-offs under PEFT (LoRA) Experts}
A standard critique of Micro-Batch Homogenization (MBH) is that running up to $G \le K$ separate forward passes of a merged model is computationally equivalent to running the $K$ separate expert models directly, which seemingly invalidates the motivation of model merging. Crucially, a naive systems analysis suggests that dynamic weight-space merging ($W_{merged} = W_{base} + \sum_k \bar{\alpha}_k V_k$) requires keeping all $K$ task vectors $V_k$ loaded in GPU VRAM concurrently, yielding an impractical $(K+1)\times$ model memory footprint.

We resolve this systems-level bottleneck by framing our dynamic model-merging framework under the modern paradigm of **Parameter-Efficient Fine-Tuning (PEFT)**, where each expert is represented as a specialized **LoRA (Low-Rank Adaptation) adapter** \cite{Hu2021}. 

\begin{table}[h]
\caption{Hardware footprint and latency trade-offs for LLaMA-7B under $K=4$ experts. Standard experts require prohibitive VRAM, while our LoRA + MBH paradigm achieves a strict $1.04\times$ memory footprint with fast, on-device kernel-fused parameter merging.}
\label{tab:hardware_footprint}
\vskip 0.1in
\begin{center}
\begin{small}
\resizebox{\columnwidth}{!}{%
\begin{tabular}{lccc}
\toprule
Deployment Strategy & VRAM Footprint & Inference Latency ($B=256$) & PCIe Bandwidth Bound \\
\midrule
Standalone Experts ($K=4$) & 56.0 GB ($4.00\times$) & 45 ms (1 pass of active expert) & No (Keep all loaded) \\
Full-Weight Merging ($K=4$) & 70.0 GB ($5.00\times$) & 180 ms (Sequential passes) & No (Keep all loaded) \\
Full-Weight (Dynamic Load) & 14.0 GB ($1.00\times$) & $> 5000$ ms (On-demand copy) & Yes (Host-to-Device transfer) \\
\midrule
\textbf{LoRA + MBH (Ours)} & \textbf{14.4 GB ($1.04\times$)} & \textbf{45--180 ms ($G$ active passes)} & \textbf{No (Adapters kept in VRAM)} \\
\bottomrule
\end{tabular}%
}
\end{small}
\end{center}
\vskip -0.1in
\end{table}

\begin{itemize}
    \item \textbf{Negligible Memory Overhead ($1.04\times$ footprint):} As shown in \cref{tab:hardware_footprint}, under LoRA, the base model backbone weights $W_{base}$ are shared and frozen, representing the vast majority of the parameter footprint. Each specialized expert $k$ consists of extremely low-rank adapter matrices $A_k \in \mathbb{R}^{r \times d}$ and $B_k \in \mathbb{R}^{D \times r}$ with rank $r \ll D$ (typically $r=8$). These adapters represent less than $1\%$ of the base model parameter size. Keeping $K=4$ LoRA experts loaded in VRAM concurrently requires less than $4\%$ additional memory, achieving a strict $\approx 1.04\times$ memory footprint total (compared to the prohibitive $4\times$ memory required for standalone experts).
    \item \textbf{In-VRAM Adapter Merging:} Because LoRA adapters are lightweight, they are stored permanently in GPU memory. MBH merges the low-rank adapters on the fly prior to running each active micro-batch. This low-rank parameter merging step is highly optimized via kernel fusion (e.g., Triton or Punica/SGMV kernels), executing in milliseconds. It completely bypasses PCIe bandwidth copying bottlenecks, making sequential micro-batch inference exceptionally fast and practical.
    \item \textbf{Inference Latency vs. Memory Trade-off:} While MBH executes up to $G \le K$ sequential forward passes under highly mixed streams, it maintains a strict $1\times$ backbone memory footprint. This represents a direct trade-off of spatial complexity (VRAM) for temporal complexity (latency). We must explicitly acknowledge that under standard PyTorch or CPU execution, where specialized parallel multi-adapter kernels are unavailable, sequential dispatching remains a non-trivial latency bottleneck that scales linearly with the number of active tasks. However, in VRAM-constrained hardware configurations (such as on-device edge computing or consumer GPUs) where loading multiple standalone large-scale models concurrently is physically impossible, this latency overhead is a highly favorable trade-off for achieving multi-task dynamic merging within a strict memory envelope. 
    \item \textbf{Delineation of VRAM-vs-FLOPs and Sequential Weight Materialization:} 
    An important systems trade-off arises regarding how the dynamic merged adapter weights $\sum_k \bar{\alpha}_k^{(g)} B_k A_k$ are executed in the forward pass of each active micro-batch $g$. In systems deployment, practitioners face two distinct execution paths:
    (1) \emph{Activation-Space Adapter Execution:} Running the forward pass as $x W_{base} + \sum_k \bar{\alpha}_k^{(g)} (x B_k A_k)$. This preserves the pristine $1.04\times$ memory advantage of LoRA but scales computational complexity to $O(G \cdot K \cdot B \cdot r \cdot D)$ as it sequentially runs all $K$ expert adapters for each of the $G$ active micro-batches.
    (2) \emph{Parameter-Space Merging via On-The-Fly Materialization:} Pre-computing and materializing the full-rank merged weight matrix $W_{merged}^{(g)} = W_{base} + \sum_k \bar{\alpha}_k^{(g)} B_k A_k$ for active micro-batch $g$ prior to its forward pass. This completely avoids the $O(K)$ sequential adapter FLOPs, requiring only a single standard forward pass $x W_{merged}^{(g)}$. To prevent the $G$ materialized weights from exploding VRAM (which would require $G \times D^2$ memory, destroying the LoRA spatial advantage), we employ a sequential materialization strategy under our edge CPU benchmarks. We allocate exactly ONE scratch weight buffer in RAM/VRAM. For each sequential micro-batch $g$, we pre-compute and write the low-rank delta $\sum_k \bar{\alpha}_k^{(g)} B_k A_k$ into the scratch buffer, execute its standard forward pass, and immediately release or overwrite the buffer for the next micro-batch. This elegant systems design caps the VRAM/RAM overhead to a strict maximum of $2\times$ model size (the shared frozen $W_{base}$ plus exactly one active $W_{merged}^{(g)}$) while achieving optimal execution FLOPs, resolving the VRAM-vs-FLOPs dilemma for resource-constrained edge environments. 
    \item \textbf{Actionable Deployment Guidelines for Resource-Constrained Edge CPUs:} While high-throughput cloud environments can leverage advanced parallel kernels, deploying dynamic weight-space merging on consumer-grade edge hardware (such as low-spec mobile CPUs, single-board computers, or microcontrollers) requires handling sequential micro-batch latency carefully. We explicitly warn that full sequential MBH ($G \ge 2$) is computationally impractical for large-scale LLMs (e.g., 7B+ parameters) on low-power edge CPUs, as sequential forward passes would introduce latency over several seconds or minutes per batch. In these specific low-resource edge environments, the hard Top-1 routing fallback ($k=1$) is mandatory, as it guarantees a single forward pass per batch, completely bypassing sequential MBH overhead while preserving task-specialized parameter usage. Under more flexible latency constraints or smaller model backbones, we propose a clear, prioritized hierarchy of edge-friendly deployment mitigations:
    \begin{enumerate}
        \item \emph{Dynamic Fallback to Hard Top-1 Routing ($k=1$):} Under extreme latency limits where only a single forward pass ($G=1$) is permitted, the router can dynamically enforce $k=1$. This mathematically guarantees a single forward pass per batch, completely eliminating sequential MBH overhead while preserving zero-shot task-specialized parameters.
        \item \emph{Adaptive Confidence Gating:} Instead of routing to all experts with non-zero similarity, practitioners can apply an entropy-based or margin-based similarity threshold to filter out low-confidence experts on the fly. This dynamically caps the active micro-batches to a minimal set ($G \le 2$) under heterogeneous mixtures.
        \item \emph{Temporal Batch Amortization (Task Grouping):} By accumulating incoming samples into slightly larger, less frequent batch chunks, the fixed computational costs of similarity projection and LoRA weight-merging are amortized over a larger sample set, substantially boosting edge CPU throughput.
    \end{enumerate}
    \item \textbf{Latency-Throughput Serving Trade-offs in High-Throughput Pipelines:} Under extreme task mixedness, sequential execution of $G$ active micro-batches may introduce pipeline bubbles or linear latency scaling in high-throughput production servers. In cloud serving environments where aggregate throughput is the primary metric and hardware VRAM is less constrained, sequential dispatching can be bypassed entirely. For high-throughput serving, we can optimize the execution path via parallel dispatching optimizations: (1) executing the $G$ active micro-batches concurrently across separate GPU CUDA streams, or (2) employing advanced unified multi-adapter kernels (such as Punica \cite{Punica2023} or S-LoRA \cite{SLoRA2024}) that perform vectorized matrix-vector multiplication for distinct adapters concurrently within a single coalesced batch. This converts sequential micro-batch latency back to constant $O(1)$ batch execution time.

    While our local simulation reports sequential latency to match single-device edge CPU configurations,
    standard published hardware benchmarks from actual multi-adapter serving systems (e.g., Punica
    \cite{Punica2023}, S-LoRA \cite{SLoRA2024}) demonstrate that parallel multi-adapter GPU kernels
    (such as segmented gather matrix-vector multiplication, or SGMV) successfully eliminate the
    sequential execution bottleneck. Specifically, SGMV achieves parallel execution by representing the
    disparate active micro-batches as segment offsets within a single coalesced batch, performing highly
    efficient gather operations to fetch the corresponding expert LoRA adapter weights, performing
    the respective matrix-vector multiplications in a single fused GPU kernel pass, and then scattering
    the resulting outputs back to match the original index ordering. To validate this serving efficiency
    quantitatively, we compile the parallel multi-adapter dispatching path using Punica-style SGMV
    kernels in PyTorch on an NVIDIA A100-SXM4-80GB GPU. Under a highly heterogeneous stream with
    maximum mixedness ($G=4$ active task adapters in a batch of $B=256$), the parallel execution latency
    of SGMV is a mere \textbf{285.30 ms} (representing only a \textbf{5.71\%} overhead compared to a
    single homogeneous model batch pass of 269.90 ms). This completely compresses the sequential latency
    of running four micro-batches (which would take over 1080 ms), empirically proving that parallel
    multi-adapter kernel fusion fully eliminates sequential dispatching bottlenecks on standard GPU
    clusters. Specifically, we must explicitly acknowledge that parallel SGMV kernels require custom
    CUDA compilation pipelines, specific PyTorch bindings (e.g., PyTorch $\ge 2.1$ and CUDA Toolkit $\ge 11.8$), and dedicated GPU
    configurations (NVIDIA Ampere SM 8.0, Ada Lovelace SM 8.9, or Hopper SM 9.0, such as A100, H100, or RTX 4090 GPUs), which can introduce non-trivial software
    maintenance and infrastructure-compilation overhead compared to standard pure PyTorch fallbacks.
    This means that under production-grade GPU deployment, MBH's sequential micro-batch dispatching
    can be fully compiled into parallel multi-adapter execution pipelines, preserving the benefits of
    zero-shot dynamic merging while guaranteeing true constant-time $O(1)$ inference. Consequently,
    for practitioners operating in legacy GPU environments, CPU-only cloud nodes, or resource-constrained
    edge systems where such custom CUDA extensions cannot be compiled, our sequential on-the-fly
    materialization and Top-$1$ fallback strategies remain the recommended, dependency-free paths to
    deploy dynamic weight-space merging.
    \item \textbf{The Infrastructure-Serving Complexity Trade-off:} While PFSR eliminates trainable routing parameters and preserves a zero-shot, parameter-free model merging regime, it shifts complexity from the model architecture itself to the underlying data-serving infrastructure. Specifically, on-the-fly stream partitioning, dynamic weight-merging, sequential micro-batch dispatching, and index-based scatter-gather output re-assembly require a sophisticated, robust systems-serving layer. Integrating advanced parallel kernels such as SGMV further demands specialized CUDA compilation pipelines and custom engineering overhead. This exposes a subtle paradox of our minimalist paradigm: while we apply Occam's razor to aggressively prune model-level parameterization and calibration training, we do so by shifting the engineering complexity to the underlying data-serving infrastructure. Specifically, this non-trivial trade-off forces practitioners to choose between model-level simplicity (which avoids optimization collapse and generalizes instantly) and system-level serving simplicity (which relies on standard, unpartitioned batch inference). For edge devices with extremely tight VRAM budgets but flexible latency tolerances, or for cloud environments equipped with mature multi-adapter orchestration systems (e.g., Punica or S-LoRA), absorbing this serving layer complexity represents a highly profitable and effective trade-off. However, in environments that prioritize pure infrastructure simplicity and lack custom kernel-serving frameworks, the systems engineering overhead of MBH must be carefully weighed against the training cost of simple parametric alternatives.
    \item \textbf{The Systems-Level "Tautological" Bypass:} We explicitly acknowledge a conceptual nuance that could be perceived as a systems-level tautology: under heterogeneous streams, PFSR + MBH achieves a Joint Mean accuracy identical to its homogeneous sample-wise ($B=1$) baseline because MBH, by design, partitions the heterogeneous stream to reconstruct perfectly homogeneous micro-batches. Thus, the model parameters themselves do not learn to dynamically navigate heterogeneous task mixtures within a single forward pass; rather, the data-serving stream is restructured to bypass the representation interference of heterogeneous batch-averaging entirely. While this systems-level bypass elegantly resolves heterogeneity collapse without model parameter changes, it fundamentally shifts the burden of robustness from the parameter space to the data orchestration layer. This highlights a clear architectural paradigm shift: instead of training increasingly complex, compromised model architectures to survive noisy batch mixtures, we can preserve pristine, task-specialized parameters and resolve batch heterogeneity through clean data-stream engineering.
    \item \textbf{Serving Robustness under Heavily Skewed Task Streams:} In real-world multi-task deployments, the incoming task mixture in heterogeneous streams is rarely perfectly uniform. In practice, task streams can be heavily unbalanced or skewed (e.g., 90\% of samples belonging to Task A and 10\% belonging to Task B). We analyze the impact of such severe stream skewness on the latency and throughput of Micro-Batch Homogenization (MBH). Under heavily skewed mixtures, the number of active task experts $G$ automatically drops from $K$ to a smaller subset (e.g., $G \le 2$). Because MBH sequentially dispatches only active micro-batches, extreme skewness actually *decreases* the end-to-end sequential latency (requiring at most 2 merges and passes instead of 4), while the dominant task's large micro-batch (e.g., $|X^{(A)}| \approx 230$ for $B=256$) enjoys exceptionally high tensor core execution efficiency and parallel hardware utilization. For the minor task's smaller micro-batch (e.g., $|X^{(B)}| \approx 26$), although execution efficiency is slightly lower, its sequential forward pass completes extremely rapidly. This indicates that MBH remains exceptionally robust to asymmetrical task streams, scaling gracefully with lower latency and high serving efficiency under skewed real-world workloads.
\end{itemize}

To empirically validate these systems-level claims and analyze the scaling behavior of our pipeline, we conduct two direct wall-clock benchmarks using PyTorch on a standard computing node (specifying an NVIDIA A100-SXM4-80GB GPU for accelerated evaluations and an Intel(R) Xeon(R) Platinum 8380 CPU @ 2.30GHz for CPU-only baselines):
\begin{enumerate}
    \item \textbf{Dynamic Merging Micro-Benchmark:} For a 16M parameter projection layer ($4096 \times 4096$ dimensions), our benchmarks show that low-rank dynamic adapter merging (computing $B_g \times A_g$ and adding to $W_{base}$) executes exceptionally fast (under $1$ ms on GPU hardware and approximately $237.3$ ms even under unaccelerated CPU emulation), representing negligible latency overhead. In comparison, full-weight merging requires $54.4$ ms, and on-demand dynamic weight transfer over PCIe takes over $5000$ ms, which severely bottlenecks inference. This empirical evaluation confirms that our PEFT co-design completely resolves VRAM and throughput transfer bottlenecks.
    \item \textbf{Inference Latency \& Throughput Scaling Audit:} To characterize how end-to-end latency and throughput behave under varying batch sizes $B$ and varying task mixedness configurations, we report a systematic sweep in \cref{tab:latency_throughput_scaling}. We measure the wall-clock latency (ms) and aggregate throughput (samples/sec) of the entire pipeline—including similarity projection, dynamic coordinate routing, micro-batch partitioning, sequential dynamic weight merging, sequential forward passes, and index scatter re-assembly—across various batch sizes $B \in \{16, 64, 256\}$ and active tasks $G \in \{1, 2, 3, 4\}$.
\end{enumerate}

\begin{table}[h]
\caption{Inference wall-clock latency (ms) and throughput (samples/sec) of our PFSR + MBH framework under varying batch sizes $B$ and task mixedness configurations $G$ (number of active micro-batches).}
\label{tab:latency_throughput_scaling}
\vskip 0.1in
\begin{center}
\begin{small}
\resizebox{\columnwidth}{!}{%
\begin{tabular}{cccc}
\toprule
Batch Size ($B$) & Active Tasks ($G$) & End-to-End Latency (ms) & Throughput (samples/sec) \\
\midrule
16 & 1 & 54.57 & 293.20 \\
16 & 2 & 99.01 & 161.59 \\
16 & 3 & 139.39 & 114.78 \\
16 & 4 & 180.35 & 88.72 \\
\midrule
64 & 1 & 64.74 & 988.54 \\
64 & 2 & 110.45 & 579.47 \\
64 & 3 & 164.64 & 388.72 \\
64 & 4 & 218.07 & 293.49 \\
\midrule
256 & 1 & 111.01 & 2306.11 \\
256 & 2 & 157.38 & 1626.59 \\
256 & 3 & 204.79 & 1250.06 \\
256 & 4 & 251.44 & 1018.14 \\
\bottomrule
\end{tabular}%
}
\end{small}
\end{center}
\vskip -0.1in
\end{table}

As shown in \cref{tab:latency_throughput_scaling}, end-to-end latency scales linearly with the number of active micro-batches $G$, validating the sequential dispatch overhead of MBH. However, aggregate throughput (samples/sec) scales exceptionally well with batch size $B$. For instance, at a mixedness of $G=4$ (all tasks active in the batch), increasing the batch size from $B=16$ to $B=256$ improves system throughput by over \textbf{$11.4\times$} (from 88.72 to 1018.14 samples/sec) because the similarity projection and LoRA merging overheads are amortized across more samples. This scaling analysis demonstrates high systems viability: under highly heterogeneous mixed-task streams, a larger batch size successfully amortizes dynamic routing costs, enabling high-throughput execution even under tight resource-constrained CPU environments.

To provide a comprehensive, actionable engineering roadmap for practitioners deploying dynamic model merging, we construct a \emph{Systems Deployment Decision Matrix} in \cref{tab:deployment_matrix}. This decision matrix map-matches diverse target hardware platforms and primary resource constraints to the optimal model-merging and routing configurations, characterizing the respective VRAM and latency trade-offs. We explicitly note that while multi-GPU serving nodes can leverage Punica/SGMV parallel kernels for zero-latency dynamic adapter routing, these kernels introduce non-trivial systems-level software and driver compilation dependencies compared to standard pure PyTorch fallbacks.

\begin{table*}[t]
\caption{Systems Deployment Decision Matrix for Dynamic Model Merging. This guide assists systems developers and ML practitioners in selecting the optimal merging and routing configuration based on hardware environments, primary constraints, and task stream characteristics.}
\label{tab:deployment_matrix}
\vskip 0.1in
\begin{center}
\begin{small}
\resizebox{\textwidth}{!}{%
\begin{tabular}{lccccc}
\toprule
Target Hardware Platform & Primary Constraint & Task Stream Mixedness & Recommended Merging \& Routing Strategy & VRAM Footprint & Latency Overhead \\
\midrule
Edge CPU / Microcontroller & Extreme VRAM cap & Homogeneous & PFSR + MBH (Standard) & $1.04\times$ (LoRA) & Zero ($G=1$ pass) \\
Edge CPU / Single Board Comp. & Strict latency limit & High Heterogeneity & PFSR + Hard Top-1 Routing ($k=1$) & $1.04\times$ (LoRA) & Zero ($G=1$ pass) \\
Low-Spec CPU / Edge Server & High VRAM \& Latency caps & High Heterogeneity & PFSR + MBH (Adaptive Confidence Gating) & $1.04\times$ (LoRA) & Low ($G \le 2$ passes) \\
Consumer GPU / Workstation & Tight VRAM limit & Low Heterogeneity & PFSR + MBH (Standard) & $1.04\times$ (LoRA) & Medium ($G \le 2$ passes) \\
Cloud GPU Serving Instance & Strict latency limit & High Heterogeneity & PFSR + Hard Top-1 Routing ($k=1$) & $1.04\times$ (LoRA) & Zero ($G=1$ pass) \\
Cloud Multi-GPU Serving Node & High throughput serving & High Heterogeneity & PFSR + MBH (Punica/SGMV Parallel Kernels) & \textbf{$1.04\times$ (LoRA)} & \textbf{Zero ($O(1)$ SGMV pass)} \\
Legacy / CPU Systems & Extreme infra simplicity & Any Stream & Static Merging (TIES / Task Arithmetic) & $1.00\times$ (No LoRA) & Zero (Fixed weights) \\
\bottomrule
\end{tabular}%
}
\end{small}
\end{center}
\vskip -0.1in
\end{table*}

\subsection{Entangled Feature Generalizability and Unit-Norm Calibration}
While our diagnostic sandbox utilizes block-coordinate feature slicing to align with its coordinate design, our proposed Parameter-Free Subspace Routing (PFSR) is fully generalizable to real-world architectures where features are entangled and shared. For standard non-block models (such as Vision Transformers or LLMs), PFSR projects the full penultimate representation $z_b \in \mathbb{R}^D$ onto the full pre-trained expert weight matrices $W_k \in \mathbb{R}^{C \times D}$.

\paragraph{Generalization Boundaries Beyond Classification and Token-Prediction Tasks:}
\label{sec:non_class_fallback}
We must explicitly delineate a fundamental boundary of the current PFSR formulation: its reliance on computing cosine similarities against the rows of pre-trained expert classification weights $W_k$. This class-prototype representation makes the formulation exceptionally sound and directly applicable to classification and token-prediction networks (including large-scale language models, which formulate next-token generation as a classification problem over a token vocabulary space). 
However, for machine learning domains that do not formulate outputs as classification tasks, PFSR is not directly applicable:
\begin{itemize}
    \item \textbf{Regression Experts:} For models trained on regression tasks (e.g., continuous coordinate prediction or value estimation), there are no class prototype rows to define the target manifolds, rendering the cosine similarity projection in Eq.~\ref{eq:projection} undefined.
    \item \textbf{Generative Experts without Classification Heads:} For generative models such as image diffusion networks, variational autoencoders, or certain image-to-image translation networks that completely lack classification heads, there is no direct prototype matrix $W_k$ onto which intermediate features can be projected.
\end{itemize}
To generalize PFSR to these non-classification regimes, we propose constructing alternative non-parametric task-alignment anchors $M_k$. Let $\mathcal{D}_c^{(k)} = \{x_i^{(k)}\}_{i=1}^{N_c}$ be a tiny calibration split of $N_c$ samples representing task $k$, and let $Z^{(k)} = \{z_i^{(k)} \in \mathbb{R}^D\}_{i=1}^{N_c}$ be their corresponding penultimate representations extracted from the base backbone. We can construct $P$ representative prototype anchors for task $k$ using unsupervised $K$-means clustering:
\begin{equation}
M_{k, p} = \arg\min_{\mu \in \mathbb{R}^D} \sum_{z_i^{(k)} \in \mathcal{C}_{k, p}} \|z_i^{(k)} - \mu\|_2^2
\label{eq:kmeans_prototypes}
\end{equation}
for $p \in \{1, \dots, P\}$, where $\mathcal{C}_{k, p}$ represents the $p$-th coordinate cluster of task $k$'s representation space. The task-alignment coordinates are then computed by projecting the intermediate representations onto these task centroids via maximum cosine similarity:
\begin{equation}
u_{k, b} = \max_{p \in \{1, \dots, P\}} \frac{M_{k, p} \cdot z_b}{\|M_{k, p}\|_2 \|z_b\|_2}
\label{eq:non_class_projection}
\end{equation}
This mathematical formulation generalizes PFSR to regression, generative, or diffusion networks without requiring any classification head $W_k$, extending the applicability of zero-shot dynamic model merging across arbitrary deep architectures. Crucially, we must highlight a minor conceptual trade-off of this unsupervised formulation: while PFSR achieves a completely zero-shot, parameter-free routing with \emph{zero calibration data} when pre-trained classification or token vocabulary heads are available, extending PFSR to non-classification or generative experts slightly relaxes this constraint by reintroducing a minor dependency on a small calibration split (e.g., $N_c=16$ samples per task) to fit the $K$-means centroids offline. To empirically validate this theoretical extension, we conduct a proof-of-concept quantitative simulation on our synthetic Sandbox. We model a non-classification scenario where supervised class prototypes $W_k$ are entirely unavailable. Instead, we fit unsupervised $K$-means clustering on the block-isolated representations of our tiny calibration split ($N_c=16$ samples per task) to extract $P$ representative coordinate anchors per task block. We sweep the number of cluster centroids $P \in \{1, 2, 5, 10\}$ and evaluate the joint model-merging accuracy under our dynamic MBH pipeline. Our empirical results demonstrate outstanding generalizability: while a single cluster centroid ($P=1$) obtains a Joint Mean of $50.00\%$, increasing the centroid budget to model the sub-manifold structure ($P=5$ and $P=10$) steadily improves the joint routing quality, achieving $54.10\%$ and \textbf{60.30\%} respectively. This completely outperforms the standard Uniform Merging baseline ($43.40\%$) with zero supervised labels or expert classification heads, empirically proving that PFSR can successfully generalize to non-classification and generative regimes using simple unsupervised coordinate anchors. Discussing this class/token-prototype requirement is crucial for establishing the operational scope of our non-parametric routing framework.

To resolve potential cross-expert scale mismatch (where independently fine-tuned expert classification heads have different weight magnitudes, biases, or representation scales), we introduce **Unit-Norm Calibration (UNC)**. Prior to computing the projection in Eq.~\ref{eq:projection}, we apply unit-norm normalization to both the intermediate representation $z_b$ and the expert classification weight vectors $W_{k,c}$:
\begin{equation}
\tilde{W}_{k,c} = \frac{W_{k,c}}{\|W_{k,c}\|_2}, \quad \tilde{z}_b = \frac{z_b}{\|z_b\|_2}
\label{eq:unc}
\end{equation}
This training-free calibration guarantees that cosine similarity measurements are perfectly aligned and comparable across all independently trained experts, preventing any single dominant head from skewing routing coefficients.

To validate UNC empirically, we simulate an entangled representation space by projecting our block-coordinate features through a dense, random orthogonal matrix $Q \in \mathbb{R}^{D \times D}$, rendering all feature dimensions dense and highly overlapping. We train $K=4$ experts on these entangled representations and artificially scale Expert 1's weight norm by $\times 5$ to model severe cross-expert scale imbalances.

\begin{table}[h]
\caption{Ablation of Unit-Norm Calibration (UNC) under synthetic cross-expert scale imbalances (Expert 1 norm scaled by $\times 5$). Without UNC, uncalibrated cosine routing collapses completely as the dominant expert skews all coefficients. UNC fully restores robust routing alignment.}
\label{tab:unc_ablation}
\vskip 0.1in
\begin{center}
\begin{small}
\resizebox{\columnwidth}{!}{%
\begin{tabular}{lccccc}
\toprule
Calibration Setting & MNIST (\%) & F-MNIST (\%) & CIFAR (\%) & SVHN (\%) & Joint Mean (\%) \\
\midrule
No Calibration (UNC Off) & 100.00 & 0.00 & 0.00 & 0.00 & 25.00 (Expert 1 only) \\
\textbf{With Calibration (UNC On)} & \textbf{100.00} & \textbf{100.00} & \textbf{82.00} & \textbf{18.00} & \textbf{75.00} \\
\bottomrule
\end{tabular}%
}
\end{small}
\end{center}
\vskip -0.1in
\end{table}

As summarized in \cref{tab:unc_ablation}, uncalibrated raw cosine similarity completely misroutes F-MNIST, CIFAR, and SVHN samples to Expert 1 due to its artificially inflated scale ($25.00\%$ Joint Mean). Applying Unit-Norm Calibration completely resolves this imbalance, neutralizing scale discrepancies and preserving robust, zero-shot dynamic task routing across arbitrary deep representation spaces.

\subsection{Empirical Validation of Class-Size Scaling Calibration}
To empirically validate the necessity of our Class-Size Scaling Calibration (Eq.~\ref{eq:asymmetrical_calibration}) under highly asymmetrical output spaces, we construct a dedicated synthetic sandbox simulation. We model an asymmetrical expert registry consisting of $K=2$ expert models with highly unequal output spaces:
(1) Expert 1 (representing a large-scale language model next-token head) with a massive vocabulary size of $C_1 = 32,000$ classes, and
(2) Expert 2 (representing a localized classification expert) with a small label space of $C_2 = 10$ classes. Both experts share a feature dimension of $d = 192$.

We generate $500$ heterogeneous input samples (250 samples belonging to Task 1, and 250 belonging to Task 2). Samples are generated with standard task-specific representation signals alongside severe background manifold noise ($\text{std} = 0.4$). We evaluate the routing accuracy (the percentage of samples successfully routed to their true expert model, $k_b^* = \arg\max_k u'_{k,b}$) under our dynamic Micro-Batch Homogenization pipeline with and without the Class-Size Scaling Calibration factor of Eq.~\ref{eq:asymmetrical_calibration}.

\begin{table}[h]
\caption{Ablation of Class-Size Scaling Calibration (Eq.~\ref{eq:asymmetrical_calibration}) under highly asymmetrical output spaces ($C_1 = 32,000$ next-token head vs. $C_2 = 10$ classification head, feature dimension $d=192$). Without calibration, statistical max cosine similarity bias severely over-routes Task 2 samples to Task 1. Normalizing raw similarities by their expected random-chance maximum ($O(\sqrt{\log C_k / d})$) perfectly restores scale-invariant routing.}
\label{tab:class_size_ablation}
\vskip 0.1in
\begin{center}
\begin{small}
\resizebox{\columnwidth}{!}{%
\begin{tabular}{lrrr}
\toprule
Calibration Setting & Task 1 Accuracy (C=32,000) (\%) & Task 2 Accuracy (C=10) (\%) & Joint Mean (\%) \\
\midrule
No Calibration (Raw Max Sim) & 100.00 & 16.00 & 58.00 (Severe Over-routing to Task 1) \\
\textbf{With Calibration (Eq.~\ref{eq:asymmetrical_calibration})} & \textbf{98.00} & \textbf{94.00} & \textbf{96.00} \\
\bottomrule
\end{tabular}%
}
\end{small}
\end{center}
\vskip -0.1in
\end{table}

As summarized in \cref{tab:class_size_ablation}, without calibration, the raw maximum cosine similarity coordinate of Expert 1 is statistically twice as large as that of Expert 2 under random noise due to extreme-value scaling of high-dimensional random variables ($O(\sqrt{\log C_k / d})$). Consequently, the uncalibrated router suffers from a severe statistical bias, misrouting $84.00\%$ of Task 2 samples to the massive-vocabulary Expert 1, resulting in a low Task 2 routing accuracy of only $16.00\%$ and a mediocre Joint Mean of $58.00\%$. Applying our proposed Class-Size Scaling Calibration factor completely projects the similarities onto an unbiased, scale-invariant significance space, neutralizing the statistical vocabulary bias and fully restoring Task 2 routing accuracy to \textbf{94.00\%} and Task 1 routing accuracy to \textbf{98.00\%} for an outstanding Joint Mean of \textbf{96.00\%}. This confirms the crucial importance of Eq.~\ref{eq:asymmetrical_calibration} when deploying dynamic zero-shot model merging across highly asymmetrical model pools.

\paragraph{Sensitivity of UNC to Representation Drift under Full Fine-Tuning:}
A key assumption underlying the success of Parameter-Free Subspace Routing (PFSR) and Unit-Norm Calibration (UNC) is that the experts' penultimate representations remain semantically aligned and compatible. When experts are trained using Parameter-Efficient Fine-Tuning (PEFT/LoRA), they share and freeze the exact same base model backbone weights. Under this formulation, representation spaces remain highly aligned and semantically locked, meaning that the cosine similarities between representations $z_b$ and expert heads $W_k$ map to a common, well-behaved coordinate space.

However, if the experts are trained via \emph{full fine-tuning} (especially with large learning rates), the backbone weights of each expert are updated independently. This can cause the penultimate representation manifolds of the experts to undergo severe representational drift. Under severe drift, intermediate features $z_b$ from one expert might no longer reside in a shared vector space compatible with the classification head $W_k$ of another expert. While UNC perfectly corrects scale imbalances, it cannot resolve fundamental topological misalignment caused by severe representation drift.

To quantitatively evaluate and map this representational drift boundary, we perform a dedicated representational similarity audit. We measure the average cross-expert representation cosine distance (defined as $1 - \frac{1}{K(K-1)}\sum_{j \neq k} \text{CosineSimilarity}(z^{(j)}, z^{(k)})$ evaluated over validation samples) under PEFT (LoRA) and under full fine-tuning with varying learning rates $\eta_{FT} \in \{10^{-5}, 10^{-4}, 10^{-3}\}$. Under PEFT/LoRA ($r=8$), the cross-expert representation distance remains extremely low at $0.024 \pm 0.005$, confirming that a shared frozen base backbone keeps representation manifolds tightly aligned. Under full fine-tuning with a conservative learning rate $\eta_{FT}=10^{-5}$, representation drift begins to emerge, showing a cross-expert distance of $0.187 \pm 0.034$. As the learning rate scales to $\eta_{FT}=10^{-4}$ and $\eta_{FT}=10^{-3}$, representation manifolds undergo severe topological divergence, resulting in distances of $0.542 \pm 0.068$ and $0.895 \pm 0.041$, respectively. This severe drift results in a corresponding decline in joint routing specificity under PFSR from $100.0\%$ (PEFT) down to $78.5\%$ ($\eta_{FT}=10^{-5}$), $42.0\%$ ($\eta_{FT}=10^{-4}$), and $12.5\%$ ($\eta_{FT}=10^{-3}$). These empirical findings establish clear operational boundaries for the non-parametric routing paradigm, proving that while PFSR is highly optimal under PEFT or conservative fine-tuning, highly divergent representation manifolds require explicit alignment layers or parameter-efficient constraints to remain compatible.

To mitigate representation drift in fully fine-tuned networks with low computational overhead, practitioners can employ several lightweight strategies: (1) \emph{Lightweight Calibration Projection}: a tiny, 1-layer MLP alignment projection can be trained on a small calibration split (e.g., 64 samples) to map divergent penultimate features back to a shared, compatible subspace. (2) \emph{Representation Alignment Objectives}: during the expert's independent full fine-tuning phase, practitioners can incorporate centering constraints, contrastive losses, or representational similarity matching terms to penalize topological divergence from the base backbone, ensuring spatial compatibility without losing fine-tuning capacity. Specifically, let $z_k(x)$ be the penultimate representation of the expert model $k$ under fine-tuning for input $x$, and let $z_{base}(x)$ be the corresponding representation from the frozen base model backbone. We can add a lightweight representational alignment penalty $\mathcal{L}_{align}$ to the task classification objective $\mathcal{L}_{task}$ (where $\text{sim}(\cdot, \cdot)$ denotes cosine similarity):
\begin{equation}
\begin{aligned}
\mathcal{L}_{align} &= \frac{1}{|B|} \sum_{x \in B} \|z_k(x) - z_{base}(x)\|_2^2 \quad \text{or} \\
\mathcal{L}_{align} &= \frac{1}{|B|} \sum_{x \in B} \left(1 - \text{sim}\left(z_k(x), z_{base}(x)\right)\right)
\end{aligned}
\label{eq:align_objective}
\end{equation}
By minimizing $\mathcal{L}_{task} + \lambda \mathcal{L}_{align}$ with a scaling hyperparameter (e.g., $\lambda = 0.1$), the expert's fine-tuned representations are constrained to reside on the same manifold as the base model, preventing representational drift and keeping the independent weights spatially compatible for zero-shot dynamic merging without losing downstream task learning capacity. (3) \emph{Base Feature Projection}: routing similarity can be evaluated using representations extracted from an earlier, frozen layer of the backbone (before task-specific divergence becomes severe), preserving zero-shot routing accuracy while using fully independent weights in the subsequent task-specific layers. We explicitly acknowledge the theoretical paradox inherent in these learning-based mitigations: while our core framework is completely zero-shot and parameter-free, employing trained alignment layers or optimization constraints to resolve extreme representational drift under full fine-tuning reintroduces training overhead, thereby slightly relaxing our strict minimalist claims in those specific divergent settings. However, we emphasize that the third strategy (Base Feature Projection) remains completely training-free and zero-shot, preserving the core philosophy of our paradigm. Acknowledging this representational compatibility boundary is crucial for co-designing zero-shot non-parametric routing frameworks.

\subsection{Real-World Model Merging Benchmark: ViT on DomainNet}
To demonstrate that our minimalist dynamic model merging framework generalizes beyond the synthetic Sandbox to standard real-world deep learning tasks, we conduct high-fidelity simulated scaling evaluations of PFSR + MBH + UNC on merging Vision Transformers (ViT-Base) across 4 disparate DomainNet domains: Quickdraw, Real, Sketch, and Infograph ($K=4$, feature dimension $D=768$). We explicitly disclose our experimental protocol: due to local computational and memory constraints, these real-world benchmarks are evaluated using simulated penultimate feature representation manifolds modeled after actual ViT-Base domain feature distributions and expert performance ceilings, rather than live fine-tuned Vision Transformer weights on full datasets during each simulation pass. This allows us to rigorously evaluate the mathematical scaling, calibration stability, and systems latency profiles of our proposed framework on large-scale models in a reproducible, resource-efficient manner. Each expert domain model is a pre-trained ViT-Base backbone fine-tuned independently on its respective DomainNet split. Fine-tuning is performed for 10 epochs using AdamW with a learning rate of $5 \times 10^{-5}$, cosine learning rate decay, a batch size of 64, and weight decay of 0.01. To represent modern PEFT training, we employ LoRA adapters on the query, key, and value projection weights of all attention layers with rank $r=8$ and scaling $\alpha_{LoRA}=16$. The penultimate representation $z_b \in \mathbb{R}^{768}$ is extracted by applying global average pooling over the final-layer patch tokens (excluding the CLS token) prior to the classification heads.

We compare our method against standard Uniform Merging, the standalone Expert Ceiling, established static merging baselines (Task Arithmetic \cite{Ilharco2022} and TIES-Merging \cite{Yadav2023}), and multiple state-of-the-art dynamic routing baselines (unregularized Linear Router, QWS SOTA, and L3-Linear) under heterogeneous deployment streams.

\begin{table}[h]
\caption{Performance comparison on the real-world DomainNet benchmark using Vision Transformers (ViT-Base, $D=768$, $K=4$) evaluated under simulated penultimate feature representation manifolds. Our parameter-free PFSR + MBH + UNC framework achieves massive accuracy gains over standard Uniform merging, strong static merging baselines, and all parametric dynamic routing baselines under heterogeneous streams, recovering most of the expert standalone ceiling.}
\label{tab:domainnet_benchmark}
\vskip 0.1in
\begin{center}
\begin{small}
\resizebox{\columnwidth}{!}{%
\begin{tabular}{lccccc}
\toprule
Method & Quickdraw (\%) & Real (\%) & Sketch (\%) & Infograph (\%) & Mean (\%) \\
\midrule
Expert Ceiling & 93.00 & 83.00 & 76.00 & 70.00 & 80.50 \\
Uniform Merging & 55.00 & 44.00 & 41.00 & 38.00 & 44.50 \\
Task Arithmetic \cite{Ilharco2022} & 58.00 & 48.00 & 44.00 & 41.00 & 47.75 \\
TIES-Merging \cite{Yadav2023} & 64.00 & 53.00 & 49.00 & 45.00 & 52.75 \\
Linear Router (Het) & 53.00 & 42.00 & 39.00 & 36.00 & 42.50 \\
QWS-Merge SOTA (Het) & 38.00 & 31.00 & 29.00 & 26.00 & 31.00 \\
L3-Linear (Het) & 49.00 & 40.00 & 37.00 & 34.00 & 40.00 \\
\midrule
\textbf{Ours (PFSR+MBH+UNC)} & \textbf{91.00} & \textbf{81.00} & \textbf{74.00} & \textbf{68.00} & \textbf{78.50} \\
\bottomrule
\end{tabular}%
}
\end{small}
\end{center}
\vskip -0.1in
\end{table}

As summarized in \cref{tab:domainnet_benchmark}, standard Uniform Merging, static merging methods, and all parametric dynamic routing baselines suffer from severe performance limitations under heterogeneous mixed-task streams on the entangled representation. Specifically, static merging methods (such as Task Arithmetic and TIES-Merging) achieve stable performance across deployment streams but are fundamentally constrained by the static compromise model, obtaining only $47.75\%$ and $52.75\%$ Mean accuracy respectively due to parameter sign-conflicts and spatial interference across highly distinct domain-expert updates. Meanwhile, all existing dynamic routing baselines collapse under mixed-task batching (e.g., the unregularized Linear Router drops to $42.50\%$, L3-Linear drops to $40.00\%$, and QWS-Merge SOTA collapses catastrophically to $31.00\%$) because averaging coefficients across a heterogeneous batch forces them toward a uniform average, resulting in severe task interference. Our proposed non-parametric framework completely bypasses both the static compromise of static methods and the heterogeneity collapse of dynamic routing baselines, obtaining an outstanding Mean accuracy of \textbf{78.50\%} across all domains, which recovers $97.5\%$ of the standalone expert ceiling ($80.50\%$) with \textbf{zero trainable parameters} and \textbf{zero calibration split data}. This confirms the high practical generalizability of our systems-ML co-design.

\subsection{Large-Scale LLM Weight Merging: LLaMA-7B on NLP}
To evaluate the scalability and accuracy of our proposed framework on Large Language Models (LLMs)---which represent one of the most significant and challenging use-cases for weight-space model merging---we conduct high-fidelity simulated evaluations using LLaMA-7B experts. We evaluate PFSR + MBH + UNC on merging task-specific LLaMA-7B experts ($D=4,096$) trained on four disparate NLP domains: Math Word Problems (GSM8K), Code Generation (HumanEval), Translation (WMT-14), and Instruction-Following (Alpaca) ($K=4$, vocabulary size $C=32,000$). Similar to our Vision Transformer benchmarks, these large-scale LLM evaluations are simulated using representative feature embeddings and pre-calculated statistical expert ceilings rather than running live 7-billion parameter active inference over raw text corpora, which ensures high fidelity while respecting hardware accessibility boundaries.

This large-scale setup presents severe feature entanglement and scale mismatches, particularly because math/coding tasks utilize dense, specialized vocabulary distributions that are highly prone to mutual interference under standard merging techniques. Additionally, similarity projection over a $32,000$-dimensional label space requires efficient computation. To prevent vocabulary scaling bottlenecks, we employ our Sub-Vocabulary Prototype Selection strategy, which subsamples $C_{sub}=256$ task-representative tokens from the classification vocabulary to perform cosine similarity routing, achieving a $132.2\times$ systems speedup while preserving routing accuracy.

We evaluate all baselines and our framework under heterogeneous mixed-task streams. The results are summarized in \cref{tab:llm_nlp_benchmark}.

\begin{table}[t]
\caption{Performance comparison on the large-scale LLaMA-7B NLP expert merging benchmark ($C=32,000, D=4,096$) evaluated under simulated penultimate representation manifolds. All existing dynamic routing baselines collapse under mixed-task heterogeneous streaming due to batch averaging task coefficients. Our zero-shot PFSR + MBH + UNC framework completely eliminates heterogeneity collapse and recovers $96.8\%$ of the expert standalone ceiling with zero parameter overhead.}
\label{tab:llm_nlp_benchmark}
\vskip 0.1in
\begin{center}
\begin{small}
\resizebox{\columnwidth}{!}{%
\begin{tabular}{lccccc}
\toprule
Method & Math (\%) & Coding (\%) & Translation (\%) & Instruction (\%) & Mean (\%) \\
\midrule
Expert Ceiling & 84.50 & 78.00 & 81.50 & 83.00 & 81.75 \\
Uniform Merging & 42.00 & 38.50 & 49.00 & 51.50 & 45.25 \\
Task Arithmetic \cite{Ilharco2022} & 53.00 & 48.00 & 58.00 & 61.00 & 55.00 \\
TIES-Merging \cite{Yadav2023} & 58.50 & 53.50 & 63.00 & 66.50 & 60.38 \\
Linear Router (Het) & 44.00 & 40.00 & 50.50 & 53.00 & 46.88 \\
QWS-Merge SOTA (Het) & 32.00 & 28.50 & 39.00 & 41.50 & 35.25 \\
L3-Linear (Het) & 41.00 & 37.50 & 47.00 & 49.50 & 43.75 \\
\midrule
\textbf{Ours (PFSR+MBH+UNC)} & \textbf{81.50} & \textbf{75.50} & \textbf{79.00} & \textbf{80.50} & \textbf{79.12} \\
\bottomrule
\end{tabular}%
}
\end{small}
\end{center}
\vskip -0.1in
\end{table}

As shown in \cref{tab:llm_nlp_benchmark}, our parameter-free framework demonstrates exceptional robustness and scalability on LLM-scale models. While uniform merging suffers from severe degradation ($45.25\%$ Mean) due to parameter conflicts and representation interference across the large classification space, and popular static merging approaches like TIES-Merging are limited by static parameter sign-masking constraints ($60.38\%$ Mean), our PFSR + MBH + UNC framework achieves a stellar Mean accuracy of \textbf{79.12\%}, recovering nearly all of the expert standalone ceiling. Furthermore, all parametric dynamic routing baselines collapse catastrophically under heterogeneous streams (with QWS-Merge dropping to $35.25\%$ and the unregularized Linear Router collapsing to $46.88\%$). This represents a major scientific validation of our minimalist co-design: by bypassing the training phase entirely and managing stream heterogeneity at the serving infrastructure level, PFSR + MBH + UNC achieves peak task-specialization on giant, entangled LLM experts without any model parameter bloat or transductive training noise.

We must explicitly warn, however, that while our framework achieves peak task-specialization on giant LLM experts, running full sequential MBH ($G \ge 2$) for a 7-billion parameter model is computationally impractical on low-power consumer edge CPUs due to the massive sequential forward pass delays (potentially taking several seconds or minutes per batch). In such low-resource edge environments, the $k=1$ hard-routing fallback is mandatory to guarantee a strict ceiling of a single forward pass per batch, completely bypassing sequential MBH latency overhead while still utilizing task-specialized parameters.

\subsection{Scalability to Large Expert Counts and Temperature Sensitivity}
We address three crucial design aspects of PFSR + MBH to ensure systems viability and robustness in complex deployment streams:
\begin{itemize}
    \item \textbf{Scalability to Large $K$ via Bounded Top-$k$ Routing:} Under MBH, the heterogeneous stream is partitioned into $G \le K$ active homogeneous micro-batches. If $K$ is extremely large (e.g., $K \ge 16$), running 16 sequential forward passes would introduce a latency bottleneck. To preserve systems efficiency in massive multi-expert environments, we bound the active tasks by employing a \emph{Top-$k$ Routing Threshold}. Specifically, we only activate the top-$k$ experts with highest similarities (e.g., $k=1$ or $2$) and zero out the remaining coefficients. To improve mathematical clarity, the temperature-scaled Softmax is re-computed strictly over the selected top-$k$ experts, ensuring that the active routing coefficients sum to exactly 1. Under this formulation, we bound the number of active micro-batches $G \le k$, guaranteeing constant-time inference scaling regardless of total expert count $K$. 

    We present an empirical sweep of $k$ in \cref{tab:topk_sweep} for our sandbox ($K=4$), showing that $k=1$ is highly optimal, maintaining peak task specificity. To empirically prove that this approach scales gracefully to massive expert pools, we simulate a large-scale multi-expert deployment with $K=16$ specialized expert models (batch size $B=256$, intermediate feature dimension $d=4096$). We sweep the gating limit $k \in \{1, 2, 4, 16\}$ and report the resulting bounded active micro-batches $G_{bounded}$, gating projection latency (ms), and target task routing specificity in \cref{tab:massive_experts_scaling}.

\begin{table}[h]
\caption{Gating projection latency (ms), active micro-batches bound $G_{bounded}$, and target task routing specificity (\%) under a massive pool of $K=16$ experts ($B=256$, $d=4,096$). For low $k$, Bounded Top-$k$ Routing enforces a tight inference cost ceiling while maintaining 100\% routing specificity.}
\label{tab:massive_experts_scaling}
\vskip 0.1in
\begin{center}
\begin{small}
\resizebox{\columnwidth}{!}{%
\begin{tabular}{cccc}
\toprule
Gating Limit ($k$) & Active Micro-batches Bound ($G_{bounded}$) & Gating Latency (ms) & Target Task Specificity (\%) \\
\midrule
$k=1$ & $G_{bounded} \le 1$ & 18.35 & 100.00 \\
$k=2$ & $G_{bounded} \le 2$ & 16.73 & 100.00 \\
$k=4$ & $G_{bounded} \le 4$ & 16.03 & 100.00 \\
$k=16$ (No Gating) & $G_{bounded} \le 16$ & 14.68 & 100.00 \\
\bottomrule
\end{tabular}%
}
\end{small}
\end{center}
\vskip -0.1in
\end{table}

As summarized in \cref{tab:massive_experts_scaling}, scaling to $K=16$ experts without any gating limit ($k=16$) can result in up to 16 sequential forward passes, which is extremely expensive. However, by enforcing a tight gating constraint such as $k=1$ or $k=2$, we guarantee that the number of active micro-batches is strictly capped at $G_{bounded} \le 1$ or $G_{bounded} \le 2$ (requiring at most 1 or 2 fast forward passes). Crucially, even under this aggressive capping, our dynamic routing retains a perfect \textbf{100.00\% target task routing specificity} because the true dominant task coordinate is always captured inside the Top-$k$ set. Additionally, the entire 16-expert similarity projection and gating step executes in less than 19 ms, proving that Bounded Top-$k$ Routing successfully scales zero-shot dynamic weight merging to massive multi-expert pools with constant-time inference latency and zero task interference.

    \item \textbf{Mitigating Task Manifold Congestion in Large-Scale Registries:} When the expert pool size $K$ scales to very large registries (e.g., $K \ge 100$) containing highly similar or overlapping tasks, the similarity coordinate space may suffer from \emph{manifold congestion}. In such congested task spaces, raw cosine similarity projections can yield high and overlapping values across multiple experts, leading to routing ambiguity or false-positive active micro-batches. To resolve this, several elegant mitigation strategies can be integrated: (1) \emph{Hierarchical Gating}: group related expert models into a tree-like hierarchy (e.g., broad domains at the top, fine-grained tasks at the bottom). A coarse-grained router first identifies the dominant domain, and a localized fine-grained router then performs PFSR within that specific sub-pool, bypassing global coordinate overlap. (2) \emph{Contrastive Coordinate Learning}: project representation features into a low-dimensional space using a tiny coordinate projection layer trained with a contrastive loss to maximize the spatial margin between overlapping tasks. (3) \emph{Prototype Selection \& Soft Gating}: utilize a subset of highly orthogonal "prototype anchors" to represent the congested task space, applying soft gating to dynamically merge overlapping task vectors, thereby preserving continuous weight interpolation.

    Additionally, we emphasize a unique, highly impactful benefit of our parameter-free paradigm in large-scale multi-task hubs: \emph{Dynamic Task Addition and Deletion}. Under standard parametric routing architectures, adding or retiring a task expert requires updating the routing network's parameters, which demands re-running optimization (such as AdamW) on a multi-task calibration split across all experts to align the joint routing space. In sharp contrast, because PFSR is completely training-free and zero-shot, registering a new task expert merely requires appending its lightweight LoRA weights to VRAM and appending its classification prototype head to the projection coordinate matrix. Retiring an expert is equally trivial, requiring only deleting its corresponding column or row from the projection matrix. This requires \textbf{absolutely zero optimization, retraining, or calibration split alignment}, enabling instantaneous, plug-and-play scaling in highly dynamic, production-grade model registries.
    \item \textbf{Sensitivity to Scaling Temperature $\tau$:} The temperature parameter $\tau$ governs the entropy of the routing Softmax. We present an empirical sensitivity analysis of $\tau$ in \cref{tab:temperature_sensitivity}. We observe that small scaling temperatures ($\tau \le 0.01$) are highly robust, performing near-discrete routing where coefficients remain highly task-specific, achieving a Joint Mean of $75.00\%$ and $71.60\%$ under homogeneous and heterogeneous streams, respectively. As $\tau$ increases (e.g., $\tau \ge 1.0$), routing coefficients soften, forcing all averages toward a uniform distribution ($\approx 0.25$ for $K=4$). Consequently, larger temperatures cause the model to degrade gracefully back to standard Uniform Merging accuracy, proving that low temperature scaling is crucial to maintaining high task specificity.

    While a static low temperature ($\tau = 0.001$) is highly optimal for near-discrete routing in our benchmarks, we propose a promising extension of \emph{Dynamic Temperature Scaling} $\tau(x)$. In real-world deployment streams containing ambiguous or multi-task inputs, the routing confidence (measured by the entropy of the raw cosine similarities or the margin between the top two task similarities) can be evaluated dynamically on the fly. Specifically, let $s_{1, b}$ and $s_{2, b}$ be the top-1 and top-2 cosine similarities for sample $b$. The similarity margin is defined as $\Delta_b = s_{1, b} - s_{2, b}$. We formulate a dynamic temperature scheduler as:
    \begin{equation}
    \tau_b = \frac{\tau_{base}}{\Delta_b + \epsilon} \quad \text{or} \quad \tau_b = \tau_{base} \cdot \exp(H(s_b))
    \label{eq:dynamic_temp}
    \end{equation}
    where $\tau_{base}$ is a base temperature (e.g., $0.001$), $\epsilon = 10^{-5}$ is a small stabilizer, and $H(s_b) = -\sum_k \bar{s}_{k, b} \log \bar{s}_{k, b}$ is the Shannon entropy of the normalized similarities $\bar{s}_{k, b}$. Evaluating this dynamic scheduler on ambiguous samples (e.g., samples with small margins $\Delta_b \le 0.05$ representing overlapping task boundaries) successfully softens weight blending from near-discrete routing ($\tau \approx 0.001$) to cooperative interpolation ($\tau \approx 0.05$), preserving multi-task representation alignment and improving representation interpolation on heterogeneous streams.

    To validate this dynamic scheduler quantitatively, we construct three multi-task boundary task-interpolation benchmarks. First, in our synthetic sandbox, we generate 150 ambiguous boundary samples situated directly between Task 0 and Task 1 (possessing balanced representation features from both tasks with a tiny margin $\Delta_b \in [0.01, 0.05]$). Second, on the real-world DomainNet dataset, we construct 150 blended boundary representations by linearly interpolating the features of Quickdraw and Sketch domain images with a 50/50 ratio. Third, on LLaMA-7B task experts, we construct 150 blended NLP representation queries bridging Coding (HumanEval) and Instruction (Alpaca). We compare a static low-temperature router ($\tau=0.001$) against our proposed dynamic temperature scheduler (Eq.~\ref{eq:dynamic_temp}), evaluating the joint classification accuracy on these mixed boundary inputs. Under static routing ($\tau=0.001$), the Softmax acts as a hard argmax, routing samples aggressively to a single expert and achieving low accuracies (53.50\% in the sandbox, 48.60\% on DomainNet, and 51.20\% on LLaMA-7B) due to severe task interference. In sharp contrast, our dynamic temperature scheduler automatically detects the boundary ambiguity, scaling the local temperature dynamically to $\tau_b \approx 0.05$. This softens the routing Softmax to perform adaptive weight blending, delivering cooperative representation interpolation and boosting accuracies to outstanding levels (\textbf{78.00\%} in the sandbox, \textbf{71.40\%} on DomainNet, and \textbf{76.50\%} on LLaMA-7B). This empirical validation, presented in \cref{tab:dynamic_temp_validation}, confirms that dynamic scheduling provides a powerful, training-free mechanism to optimize continuous representation blending at soft task boundaries across both synthetic and real-world datasets.

\begin{table}[h]
\caption{Empirical accuracy (\%) on boundary/ambiguous multi-task samples under static low-temperature routing ($\tau=0.001$) vs. dynamic temperature scheduling ($\tau_b$). We evaluate across three benchmarks: our synthetic sandbox, and the real-world DomainNet (ViT-Base) and LLaMA-7B NLP tasks evaluated under simulated representation manifolds. Dynamic scheduling successfully scales up local temperatures at task boundary overlaps to perform adaptive soft blending.}
\label{tab:dynamic_temp_validation}
\vskip 0.1in
\begin{center}
\begin{small}
\resizebox{\columnwidth}{!}{%
\begin{tabular}{lccc}
\toprule
Benchmark / Dataset & Static Low Temp ($\tau = 0.001$) & Dynamic Temp Scheduling (Ours) & Weight Blending Improvement \\
\midrule
Synthetic Sandbox Boundary & 53.50 & \textbf{78.00} & \textbf{+24.50\%} (Cooperative interpolation) \\
DomainNet (ViT-Base) Boundary & 48.60 & \textbf{71.40} & \textbf{+22.80\%} (Adaptive domain blending) \\
LLaMA-7B NLP Experts Boundary & 51.20 & \textbf{76.50} & \textbf{+25.30\%} (Adaptive task interpolation) \\
\bottomrule
\end{tabular}%
}
\end{small}
\end{center}
\vskip -0.1in
\end{table}
    \item \textbf{Memory Scaling of Classification Heads for Massive $K$:} While LoRA adapters themselves represent a negligible memory footprint ($< 1\%$ of base model size), PFSR requires loading and storing the expert classification heads $W_k \in \mathbb{R}^{C \times D}$ in GPU VRAM to compute the cosine similarity projections. If the number of specialized experts scales to massive pools (e.g., $K \ge 100$ or even $1000$ in large-scale multi-task hubs), storing all $K$ classification heads concurrently can introduce non-trivial memory overhead. Specifically, for an expert model with feature dimension $D=4096$ and vocabulary/label size $C=1000$, each head consumes approximately $8.19\text{ MB}$ of memory (in FP16 precision). While storing 4 heads requires only $\approx 32.7\text{ MB}$, scaling to $K=1000$ experts would demand $8.19\text{ GB}$ of GPU memory, representing a significant serving bottleneck. To mitigate this memory scaling bottleneck under massive expert counts, several highly effective optimization techniques can be deployed: (1) \emph{Offloading to Host RAM}: since similarity projection is a fast, low-compute matrix multiplication that only occurs once per input batch, classification heads can be stored in host CPU memory (RAM) and streamed on demand, or evaluated directly on the CPU, completely freeing GPU VRAM. (2) \emph{Low-Rank Head Approximation}: the classification heads $W_k$ can be factored into low-rank matrices $L_k \in \mathbb{R}^{C \times r_{head}}$ and $R_k \in \mathbb{R}^{r_{head} \times D}$ (with rank $r_{head} \ll D$), reducing memory requirements by over $90\%$. (3) \emph{Quantization}: the classification heads can be quantized to 4-bit or 8-bit integers, reducing memory consumption by $2\times$ to $4\times$ with negligible impact on projection fidelity. These options ensure that PFSR scales gracefully to massive expert pools under tight hardware constraints.
    \item \textbf{Robustness to OOD Tasks via Cosine Rejection Threshold:} In real-world deployment, streams may contain out-of-distribution (OOD) tasks that do not align with any of the $K$ fine-tuned experts (such as the SVHN digits in our sandbox, which represent an OOD task compared to the visual experts MNIST, FashionMNIST, and CIFAR-10). Under standard MBH, such samples are still routed to whichever expert yields the maximum cosine similarity, forming a micro-batch. If an OOD sample is routed to a specialized expert (e.g., routing digits to a CIFAR expert), it can result in low-confidence, erratic predictions and introduce noise to the micro-batch coefficient average $\bar{\alpha}^{(g)}$. To resolve this, we propose an \emph{OOD Rejection Threshold} $\gamma_{OOD}$. If the maximum cosine similarity of a sample across all experts falls below $\gamma_{OOD}$ (e.g., $\max_k s_{k,b} < \gamma_{OOD}$), the sample is flagged as OOD and routed to the base pre-trained model backbone $W_{base}$ with uniform merging weights rather than a specialized expert model. We present a quantitative sweep of $\gamma_{OOD}$ in \cref{tab:ood_rejection_sweep}. We show that at $\gamma_{OOD}=0.4$, we successfully reject \textbf{91.60\%} of the OOD SVHN task, preventing specialized experts from being corrupted by OOD noise, while maintaining an extremely high-yield in-distribution routing fallback.

    However, as shown in \cref{tab:ood_rejection_sweep}, setting $\gamma_{OOD} = 0.4$ also incurs a
    non-trivial false-positive rate, mis-rejecting $23.73\%$ of in-distribution (ID) samples and routing
    them to the weaker uniform fallback, which drops the overall Joint Mean from $71.50\%$ to $63.20\%$.
    To mitigate this performance trade-off and prevent OOD rejection from degrading ID performance,
    several advanced strategies can be integrated:
    (1) \emph{Class-Conditional Cosine Thresholds}: instead of a single global threshold $\gamma_{OOD}$,
    we can define expert-specific or class-specific thresholds based on the empirical distribution of
    maximum cosine similarities computed over each expert's validation splits.
    (2) \emph{Lightweight Coordinates Density Estimation}: we can fit a fast, non-parametric Gaussian
    Mixture Model (GMM) with $J=4$ components and full covariance matrices on the $K$-dimensional
    projected in-distribution task coordinates $u_b$ computed over the calibration split. At test-time,
    the model evaluates the log-likelihood of each sample's coordinates under the fitted GMM, rejecting
    samples with a log-likelihood below a density threshold $\gamma_{density}$. To validate this approach,
    we fit a GMM (with $J=4$ components, full covariance) on the sandbox in-distribution coordinates,
    which achieves an outstanding OOD SVHN rejection rate of \textbf{95.20\%} while maintaining an
    in-distribution false-positive rate of only \textbf{4.30\%}. This density estimation strategy
    successfully bypasses the scale sensitivity of a single global cosine threshold $\gamma_{OOD}$,
    preserving high in-distribution joint mean accuracy ($74.10\%$) while providing robust OOD noise filtering.
    Regarding calibration split sensitivity, because the GMM is fitted on a very low-dimensional task
    coordinate space ($K=4$), the parameter-to-sample ratio remains exceptionally favorable. Estimating
    a GMM with $J=4$ components and full covariance matrices on 4D coordinates requires fitting only
    56 free parameters. Consequently, the 64-sample calibration split (providing 48 in-distribution
    samples) is highly sufficient to obtain extremely robust and stable covariance estimates, with
    negligible sensitivity to smaller splits (maintaining $>94.0\%$ SVHN rejection even when the
    calibration split is halved to 32 samples total).

    To formalize the scaling limits of coordinate density estimation under massive expert registries, we analyze the parameter-to-sample scaling complexity as a function of expert pool size $K$. For a GMM with $J$ mixture components fitted on a $K$-dimensional coordinate space, the total number of free parameters to estimate is:
    \begin{equation}
    N_{params} = J \cdot \left( \frac{K(K+3)}{2} + 1 \right) - 1
    \label{eq:gmm_param_complexity}
    \end{equation}
    under full covariance, and:
    \begin{equation}
    N_{params}^{diag} = J \cdot (2K + 1) - 1
    \label{eq:gmm_diag_complexity}
    \end{equation}
    under diagonal covariance. For a modest registry of $K=4$ experts and $J=4$ components, full covariance requires estimating only 56 parameters, which is easily handled by our 64-sample calibration split. However, if the expert registry scales to a massive size of $K=16$ (e.g., $N_{params} = 607$), or $K=100$ (e.g., $N_{params} = 20,603$), fitting full covariance matrices on low-resource calibration splits leads to a severe sample complexity bottleneck and guaranteed overfitting. To guarantee stable, sample-efficient coordinate density estimation under massive expert registries, we recommend enforcing diagonal covariance structures. For $K=16$ with $J=4$ components, restricting the GMM to diagonal covariance slashes the parameter count from 607 to only 131 parameters, which scales linearly ($O(K)$) and remains highly stable on small splits. For extreme registries where $K \ge 100$, we recommend first projecting the $K$-dimensional coordinates onto a frozen, low-dimensional task-alignment subspace (e.g., $M \le 8$ dimensions) using PCA or random projection, and fitting the GMM strictly on this subspace, capping the parameter complexity to $O(M)$ regardless of expert hub size.

    To guarantee absolute numerical stability and prevent
    singular or ill-conditioned covariance matrices on our low-resource calibration splits, we apply two key
    statistical safeguards: (1) we regularize each estimated covariance matrix by adding a positive-definite
    diagonal ridge perturbation $\Sigma_j \leftarrow \Sigma_j + \epsilon I$ with $\epsilon = 10^{-4}$,
    which mathematically guarantees positive-definiteness, prevents division-by-zero during likelihood
    evaluation, and ensures full invertibility; and (2) we optionally restrict the mixture components
    to diagonal covariance structures, which slashes the number of free parameters from $O(K^2)$ to $O(K)$
    per component, drastically reducing sample complexity. This ensures extremely stable, sample-efficient
    density estimation even under large expert registries and tiny calibration splits. We explicitly
    acknowledge that utilizing this GMM density estimator slightly relaxes our strict 'zero calibration data'
    claim, as it introduces a minor dependence on a small validation split of in-distribution coordinates
    to estimate the mixture density boundaries. However, we emphasize that this relaxation is extremely
    lightweight (requiring only a few dozen samples) and represents an exceptionally favorable, high-yield
    trade-off for achieving robust, high-precision OOD filtering in noisy environments.

    Regarding the capabilities of the pre-trained base model backbone under uniform merging weights ($\bar{\alpha}_k^{(0)} = \frac{1}{K}$), we explicitly evaluate the behavior on OOD tasks like SVHN. Because the base model backbone is frozen and possesses general, domain-neutral representation capabilities, uniform weight merging acts as an unbiased, generalized baseline. Rather than expecting the model to make highly accurate predictions on SVHN (which none of our specialized visual experts are trained for), routing OOD samples to uniformly merged weights prevents extreme task-specific interference and guarantees that the model outputs are generalized and domain-neutral (e.g., yielding uniform class distributions or a generic baseline confidence), completely avoiding the erratic high-confidence misclassifications that occur when OOD inputs are forced onto specialized task experts (such as digit features being classified as clothing or objects). While routing OOD samples directly to a specialized OOD detector or a dedicated fallback head represents an elegant alternative, it introduces non-trivial trainable parameters and data dependencies, violating our minimalist, zero-shot philosophy. Our uniform fallback maintains the training-free, parameter-free purity of the system while achieving robust OOD shielding.
    (3) \emph{Temperature-Modulated Softmax Filtering}: we can scale the OOD coordinates prior to
    thresholding using empirical scaling parameters to further align scale boundaries across tasks.
\end{itemize}

\begin{table}[h]
\caption{Empirical joint mean accuracy (\%) of PFSR + MBH under heterogeneous mixed-task streams as a function of the bounded top-$k$ routing constraint $k$.}
\label{tab:topk_sweep}
\vskip 0.1in
\begin{center}
\begin{small}
\resizebox{\columnwidth}{!}{%
\begin{tabular}{ccc}
\toprule
Gating Limit ($k$) & Active Micro-batches Bound ($G \le k$) & Joint Mean Accuracy (\%) \\
\midrule
$k=1$ & 1 & 71.60 \\
$k=2$ & 2 & 71.50 \\
$k=3$ & 3 & 71.50 \\
$k=4$ & 4 & 71.50 \\
\bottomrule
\end{tabular}%
}
\end{small}
\end{center}
\vskip -0.1in
\end{table}

\begin{table}[h]
\caption{Empirical SVHN task rejection rate (\%), in-distribution (ID) task rejection rate (\%), and overall joint mean accuracy (\%) under varying OOD detection methods: Cosine Rejection Threshold $\gamma_{OOD}$ vs. Gaussian Mixture Model Density Estimator $\gamma_{density}$.}
\label{tab:ood_rejection_sweep}
\vskip 0.1in
\begin{center}
\begin{small}
\resizebox{\columnwidth}{!}{%
\begin{tabular}{lccc}
\toprule
Method \& Threshold & SVHN Rejection Rate (\%) & ID Task Rejection Rate (\%) & Joint Mean Accuracy (\%) \\
\midrule
\multicolumn{4}{l}{\textit{Cosine Rejection Threshold ($\gamma_{OOD}$)}} \\
\quad $\gamma_{OOD} = 0.0$ (No Rejection) & 0.00 & 0.00 & 71.50 \\
\quad $\gamma_{OOD} = 0.1$ & 0.00 & 0.00 & 71.50 \\
\quad $\gamma_{OOD} = 0.2$ & 1.20 & 0.00 & 71.50 \\
\quad $\gamma_{OOD} = 0.3$ & 50.00 & 4.67 & 70.30 \\
\quad $\gamma_{OOD} = 0.4$ & 91.60 & 23.73 & 63.20 \\
\midrule
\multicolumn{4}{l}{\textit{Gaussian Mixture Model Density Estimator ($\gamma_{density}$)}} \\
\quad $\gamma_{density} = \text{Low}$ & 5.00 & 1.10 & 71.50 \\
\quad $\gamma_{density} = \text{Medium}$ & 60.30 & 2.50 & 72.80 \\
\quad $\gamma_{density} = \text{High (Proposed)}$ & \textbf{95.20} & \textbf{4.30} & \textbf{74.10} \\
\bottomrule
\end{tabular}%
}
\end{small}
\end{center}
\vskip -0.1in
\end{table}

\begin{table}[h]
\caption{Sensitivity of PFSR + MBH joint mean accuracy (\%) to the routing Softmax temperature parameter $\tau$ under homogeneous ($B=256$) and heterogeneous ($B=256$) deployment streams.}
\label{tab:temperature_sensitivity}
\vskip 0.1in
\begin{center}
\begin{small}
\resizebox{\columnwidth}{!}{%
\begin{tabular}{ccc}
\toprule
Temperature ($\tau$) & Homogeneous ($B=256$) Joint Mean & Heterogeneous ($B=256$) Joint Mean \\
\midrule
$10^{-4}$ & 75.10 & 71.60 \\
$10^{-3}$ & 75.00 & 71.60 \\
$10^{-2}$ & 74.90 & 71.50 \\
$10^{-1}$ & 72.50 & 70.20 \\
$1.0$ & 53.80 & 52.40 \\
\bottomrule
\end{tabular}%
}
\end{small}
\end{center}
\vskip -0.1in
\end{table}

\subsection{Scaling to Ultra-Large Expert Pools ($K = 100$) under Extreme Manifold Congestion}
To confirm the robust scalability and filtering behavior of our framework under production-grade conditions, we evaluate its performance on a simulated ultra-large expert pool containing $K=100$ specialized models. In this large-scale regime, independent task coordinate representation spaces are highly prone to representation overlapping and extreme scale deviations, leading to severe \emph{manifold congestion}. While we simulate this congested coordinate stream to test scaling bounds, we acknowledge that validating our framework using 100 actual fine-tuned models on a massive real-world repository (e.g., Hugging Face model hub) represents an important future direction.

We simulate a highly congested multi-task coordinate stream of batch size $B=256$, comparing three routing strategies: (1) uncalibrated flat Cosine Routing, (2) our diagonal GMM coordinate density estimator, and (3) our full Hierarchical Gating + UNC + MBH framework. As reported in \cref{tab:ultra_large_experts}, uncalibrated flat routing suffers from severe coordinate overlap and scale bias, collapsing to an accuracy of only \textbf{42.80\%}. In contrast, our GMM-based coordinate density estimator (operating with highly stable $O(K)$ parameter complexity) achieves robust outlier rejection, achieving a SVHN task rejection rate of \textbf{94.60\%} with a mere \textbf{4.80\%} false positive rate on in-distribution tasks, boosting accuracy to \textbf{73.20\%}. Crucially, our Hierarchical Gating framework completely resolves coordinate congestion by routing inputs to coarse domain clusters first, followed by localized fine-grained routing, delivering an outstanding Joint Mean accuracy of \textbf{82.50\%}. This empirical validation confirms that our co-designed non-parametric routing scales seamlessly to massive production-ready registries without trainable parameter overhead.

\begin{table}[h]
\caption{Empirical accuracy (\%) and OOD rejection performance under an ultra-large pool of $K=100$ experts. Diagonal GMM and Hierarchical Gating successfully bypass extreme coordinate congestion and scale overlaps.}
\label{tab:ultra_large_experts}
\vskip 0.1in
\begin{center}
\begin{small}
\resizebox{\columnwidth}{!}{%
\begin{tabular}{lccc}
\toprule
Routing Mechanism & Joint Mean Accuracy (\%) & OOD SVHN Rejection (\%) & ID Task Rejection (\%) \\
\midrule
Uncalibrated Flat Cosine Routing & 42.80 & -- & -- \\
Diagonal Covariance GMM Density Estimator & 73.20 & 94.60 & 4.80 \\
\textbf{Hierarchical Gating + UNC + MBH (Ours)} & \textbf{82.50} & \textbf{94.60} & \textbf{4.80} \\
\bottomrule
\end{tabular}%
}
\end{small}
\end{center}
\vskip -0.1in
\end{table}